% Siglas do TCC (ADR 0027), carregadas pelo \loadglsentries do tcc.tex.
%
% No .org, a sigla se escreve [[gls:rotulo]] ([[glspl:rotulo]] no plural,
% [[Gls:rotulo]] no começo de frase). A primeira ocorrência imprime o campo
% first (firstplural no plural), e as seguintes, o text (plural). As duas
% formas de definição do AGENTS.md:
%   - nome em português: "\emph{nome} (SIGLA), do inglês \emph{name}";
%   - nome só em inglês: "SIGLA (\emph{name})".
% Sigla corrente da área (HTTP, HTML, XML) tem first igual a text: entra na
% lista, mas não se define no texto.
%
% A description é o que sai na lista de siglas; sort, a ordem alfabética.

\newglossaryentry{api}{type=\acronymtype, name={API}, sort={API},
  text={API}, plural={APIs},
  first={API (\emph{Application Programming Interface})},
  firstplural={APIs (\emph{Application Programming Interfaces})},
  description={\emph{Application Programming Interface}}}

\newglossaryentry{bdtd}{type=\acronymtype, name={BDTD}, sort={BDTD},
  text={BDTD}, first={\emph{Biblioteca Digital Brasileira de Teses e Dissertações} (BDTD)},
  description={Biblioteca Digital Brasileira de Teses e Dissertações}}

\newglossaryentry{bom}{type=\acronymtype, name={BOM}, sort={BOM},
  text={BOM}, first={BOM (\emph{Bill of Materials})},
  description={\emph{Bill of Materials}}}

\newglossaryentry{dc}{type=\acronymtype, name={DC}, sort={DC},
  text={DC}, plural={DCs},
  first={\emph{Dimensão Cognitiva de Notações} (DC), do inglês \emph{Cognitive Dimension of Notations}},
  firstplural={\emph{Dimensões Cognitivas de Notações} (DCs), do inglês \emph{Cognitive Dimensions of Notations}},
  description={Dimensões Cognitivas de Notações, do inglês \emph{Cognitive Dimensions of Notations}}}

\newglossaryentry{dom}{type=\acronymtype, name={DOM}, sort={DOM},
  text={DOM}, first={DOM (\emph{Document Object Model})},
  description={\emph{Document Object Model}}}

\newglossaryentry{pf}{type=\acronymtype, name={PF}, sort={PF},
  text={PF}, first={\emph{programação funcional} (PF)},
  description={Programação funcional}}

\newglossaryentry{pfr}{type=\acronymtype, name={PFR}, sort={PFR},
  text={PFR}, first={\emph{programação funcional reativa} (PFR)},
  description={Programação funcional reativa}}

\newglossaryentry{poo}{type=\acronymtype, name={POO}, sort={POO},
  text={POO}, first={\emph{programação orientada a objetos} (POO)},
  description={Programação orientada a objetos}}

\newglossaryentry{pr}{type=\acronymtype, name={PR}, sort={PR},
  text={PR}, first={\emph{programação reativa} (PR)},
  description={Programação reativa}}

\newglossaryentry{xaml}{type=\acronymtype, name={XAML}, sort={XAML},
  text={XAML}, first={XAML (\emph{Extensible Application Markup Language})},
  description={\emph{Extensible Application Markup Language}}}

% Correntes da área: sem definição no texto, só na lista.
\newglossaryentry{html}{type=\acronymtype, name={HTML}, sort={HTML},
  text={HTML}, first={HTML},
  description={\emph{HyperText Markup Language}}}

\newglossaryentry{http}{type=\acronymtype, name={HTTP}, sort={HTTP},
  text={HTTP}, first={HTTP},
  description={\emph{Hypertext Transfer Protocol}}}

\newglossaryentry{ide}{type=\acronymtype, name={IDE}, sort={IDE},
  text={IDE}, first={IDE},
  description={\emph{Integrated Development Environment}}}

\newglossaryentry{xml}{type=\acronymtype, name={XML}, sort={XML},
  text={XML}, first={XML},
  description={\emph{Extensible Markup Language}}}
